\begin{abstract}
针对六区域异构 AI 任务与新能源、储能及电网协同调度问题，本文建立了以任务执行区域、开工时刻和储能动作为决策变量，以运行成本、碳排放、网络时延、新能源利用率、服务质量及峰值净购电为评价指标的多目标模型，并纳入算力、功率、时延、期限、能量平衡和 SOC 边界约束。求解阶段采用固定 Huber 滞后回归、确定性贪心任务调度、规则型 SOC 前向策略和固定种子情景复算，所得结果仅为通过硬约束审计的启发式可行方案。结果表明，$50000$ 条任务记录均通过输入校验，缺失单元格数量为 $0$，预测总体 MAE 和 RMSE 分别为 $24.9733$~GPU 和 $35.6945$~GPU。问题二所得运行成本为 $-419658147.7076$~元，碳排放为 $14.3916$~tCO$_2$，网络时延为 $5.1594$~ms，新能源利用率为 $0.6791$，迁移任务数为 $209$ 个。问题三累计充电 $2076.3155$~MWh、放电 $0$~MWh，含储能方案的运行成本增加 $304223.9291$~元，碳排放与峰值净购电变化均为 $0$。问题四联合基准的运行成本、碳排放、新能源利用率和峰值净购电分别为 $-419353923.7785$~元、$14.3916$~tCO$_2$、$0.6791$ 和 $40.4664$~MW，并实际检查了 $15$ 个情景。研究结果说明，该方法可生成满足已审计硬约束的联合调度方案，但其属于启发式复算结果，不构成 Pareto 解，且不具有最优性证书。
\end{abstract}

\noindent\textbf{关键词：}异构 AI 任务；算力调度；Huber 回归；储能协同；碳感知调度；情景分析